% !TEX root = ideas.tex
%\newcommand{\anote}[1]{\todo[inline, size=\small, author=Arkady]{#1}}
\newcommand{\tuple}[2]{\langle #1, #2 \rangle }
\newcommand{\tup}[1]{\langle #1, #1 \rangle}
\subsubsection{Server-Checked Complaints}
\paragraph{A fatal flaw?}
In the process of writing this section, I realized that my approach has a fatal flaw.  The problem is with reusing the same check poly, $cp$, for all messages.  Now, when a party $u_i$ receives any two messages, he can compute two ciphertexts with the same value of $cp(i)$.  He can then subtract these homomorphically, to get a ciphertext with a 0 in the $cp$ slot.  This enables him to change the $p(i)$ value in any subsequent complaint without violating correctness of the check poly.  

\paragraph{Original Scheme:}
One problem that we have is clients submitting invalid complaints.  Ideally, we would like to allow the Server to check whether a complaint is valid without learning what message the complaint corresponds to.

Below is a proposed solution for this, though many details in particular, the originating procedure needs to be optimized.

\paragraph{Crypto Tools:}
We use a secret-key homomorphic encryption (HE) scheme that operates on tuples of values $(x_1,x_2)$.  The HE scheme needs to be able to compute additive homomorphism and one level of multiplications.
We need the following security properties (satisfied by many HE schemes)
\begin{itemize}
\item It is hard for someone not knowing $sk$ to produce a valid ciphertext from scratch
\item When given Enc(f(x)) produced by homomorphism, it is not possible to tell what $f$ was.  A stronger form of this, known as circuit privacy, requires that one (even with $sk$) cannot tell the difference between a freshly encrypted $ct$ and one produced via homomorphism.
\end{itemize}
\paragraph{Setup:}
\begin{enumerate}
\item $S$ generates $sk$ for HE.
\item $S$ picks a random degree $t$ check poly, $cp(x)$.  Let $cp_0,\ldots,cp_t$ be the coefficients of $cp$.
\item When a user $u_i$ joins, $S$ computes $Enc_{sk}(\tup{i} ,Enc_{sk}(\tup{i^2}),\ldots,Enc_{sk}(\tup{i^t})$ and gives these to $u_i$.
\end{enumerate}

\paragraph{Originating a message: $S$ and $u_j$ where $j$ is the originator's ID}
\begin{enumerate}
\item $S$ picks a random degree $t$ poly $p_S(x)$ s.t. $p_S(0)=j$
\item $u_j$ picks a random degree $t$ poly $p_O(x)$ s.t. $p_O(0)=0$
\item $S$ and $u_j$ run a maliciously secure 2-PC where.  I am fairly confident that a generic 2-PC is not needed here, but haven't come up with something better. 
\begin{itemize}
\item $S$ inputs $sk$, $\seckeyserv$, $p_S$, $cp$
\item $u_j$ inputs $p_O$
\item The computation checks that $p_O$ is indeed a poly of degree $t$ with $p_O(0)=0$.  Then, computes $p(x)=p_S(x)+p_O(x)$ by adding the coefficients and then returns to $u_j$ encrypted coefficients of $p(x)$ together with a signature by the server authenticating them.\\
Output (to $u_j$):  \\$c = (Enc_{sk}(\tuple{p_0}{cp_0}), \ldots, Enc_{sk}(\tuple{p_t}{cp_t})), \sigma=Sig_{\seckeyserv}(c)$ 
\end{itemize}
\item $u_j$ includes $(c,\sigma)$ with every copy of the message he sends out.
\end{enumerate}

\paragraph{Complaint:}
To complain about a received message $m$, a party $u_i$ does the following:
\begin{enumerate}
\item $u_i$ first checks that $(c,\sigma)$ included in the message is a valid signature.  If not, he sends the invalid message to $S$.  \anote{May need to include the message in $\sigma$ to prevent replay attack?}
\item $u_i$ uses $Enc_{sk}(\tup{i} ,Enc_{sk}(\tup{i^2}),\ldots,Enc_{sk}(\tup{i^t})$ received from $S$ and $Enc_{sk}(\tuple{p_0}{cp_0}), \ldots, Enc_{sk}(\tuple{p_t}{cp_t})$ received in the message to compute $Enc_{sk}(\tuple{p(i)}{cp(i)})$.
\item $u_i$ sends $Enc_{sk}(\tuple{p(0)}{cp(0)})$ to $S$ who decrypts it and checks that $cp(i)$ is correct.  (Recall that $S$ knows $cp$).
\item If this check passes, $u_i$ does the PIR and write to new location trick to write this complaint to a random location (Note that $S$ will see if he write something other than what he checked in the prior step).  
\end{enumerate}
